
\documentclass{article}
\usepackage{colortbl}
\usepackage{makecell}
\usepackage{multirow}
\usepackage{supertabular}

\begin{document}

\newcounter{utterance}

\centering \large Interaction Transcript for game `wordle', experiment `medium\_frequency\_words\_no\_clue\_no\_critic', episode 7 with refineonly{-}Qwen3.5{-}4B.
\vspace{24pt}

{ \footnotesize  \setcounter{utterance}{1}
\setlength{\tabcolsep}{0pt}
\begin{supertabular}{c@{$\;$}|p{.15\linewidth}@{}p{.15\linewidth}p{.15\linewidth}p{.15\linewidth}p{.15\linewidth}p{.15\linewidth}}
    \# & \multicolumn{2}{c}{Player} && \multicolumn{2}{c}{Game Master} \\
    \hline

    \theutterance \stepcounter{utterance}  
    & & & \multicolumn{4}{p{0.6\linewidth}}{
        \cellcolor[rgb]{0.9,0.9,0.9}{
            \makecell[{{p{\linewidth}}}]{
                \texttt{\tiny{[P1$\langle$GM]}}
                \texttt{You are a language wizard who likes to guess words by using the given rules.} \\
\\ 
\texttt{Welcome to Wordle! You have six attempts to guess the target word, a valid English word of five lowercase letters (a{-}z). Please use the tags "explanation:" and "guess:" to provide a concise explanation for each guess.} \\
\\ 
\texttt{For instance, if your guess is "apple", your response should be} \\
\texttt{explanation: this is a common five{-}letter English word, and I am starting my guess with this word.} \\
\texttt{guess: apple} \\
\\ 
\texttt{After each guess, your answer will be validated, and you will receive feedback indicating which letters are correct (green), which letters are correct but in the wrong position (yellow), and which letters are incorrect (red). This feedback can be useful in determining which letters to include or exclude in your next guess.} \\
\\ 
\texttt{For example, the feedback for "apple" might be:} \\
\texttt{guess\_feedback: a<yellow> p<yellow> p<green> l<yellow> e<red>} \\
\\ 
\texttt{The explanation should contain details about how the guess\_feedback is used to arrive at a new guess.} \\
\\ 
\texttt{Do not generate any other text.} \\
\\ 
\texttt{Let's begin with your first guess.} \\
            }
        }
    }
     \\ \\

    \theutterance \stepcounter{utterance}  
    & \multicolumn{4}{p{0.6\linewidth}}{
        \cellcolor[rgb]{0.561,0.737,0.561}{
            \makecell[{{p{\linewidth}}}]{
                \texttt{\tiny{[P1$\rangle$GM]}}
                \texttt{explanation: I will start with a very common five{-}letter word that contains a variety of frequent letters to maximize information gain.} \\
\texttt{guess: apple} \\
            }
        }
    }
    & & \\ \\

    \theutterance \stepcounter{utterance}  
    & & & \multicolumn{2}{p{0.3\linewidth}}{
        \cellcolor[rgb]{0.9,0.9,0.9}{
            \makecell[{{p{\linewidth}}}]{
                \texttt{\tiny{[GM$|$GM]}}
                \texttt{attempts = 1} \\
\texttt{target\_word = owner} \\
\texttt{guess = apple} \\
\texttt{guess\_feedback = a<red> p<red> p<red> l<red> e<yellow>} \\
            }
        }
    }
    & & \\ \\

    \theutterance \stepcounter{utterance}  
    & & & \multicolumn{4}{p{0.6\linewidth}}{
        \cellcolor[rgb]{0.9,0.9,0.9}{
            \makecell[{{p{\linewidth}}}]{
                \texttt{\tiny{[P1$\langle$GM]}}
                \texttt{guess\_feedback: a<red> p<red> p<red> l<red> e<yellow>} \\
\\ 
\texttt{Provide your response only in this format.} \\
\texttt{explanation: details} \\
\texttt{guess: word} \\
            }
        }
    }
     \\ \\

    \theutterance \stepcounter{utterance}  
    & \multicolumn{4}{p{0.6\linewidth}}{
        \cellcolor[rgb]{0.561,0.737,0.561}{
            \makecell[{{p{\linewidth}}}]{
                \texttt{\tiny{[P1$\rangle$GM]}}
                \texttt{explanation: The feedback indicates that 'e' is in the word but in the last position, while 'a', 'p', and 'l' are not present. I need to find a word ending in 'e' that does not contain 'a', 'p', or 'l'. I will try a common word ending in 'e' with different vowels and consonants.} \\
\texttt{guess: crane} \\
            }
        }
    }
    & & \\ \\

    \theutterance \stepcounter{utterance}  
    & & & \multicolumn{2}{p{0.3\linewidth}}{
        \cellcolor[rgb]{0.9,0.9,0.9}{
            \makecell[{{p{\linewidth}}}]{
                \texttt{\tiny{[GM$|$GM]}}
                \texttt{attempts = 2} \\
\texttt{target\_word = owner} \\
\texttt{guess = crane} \\
\texttt{guess\_feedback = c<red> r<yellow> a<red> n<yellow> e<yellow>} \\
            }
        }
    }
    & & \\ \\

    \theutterance \stepcounter{utterance}  
    & & & \multicolumn{4}{p{0.6\linewidth}}{
        \cellcolor[rgb]{0.9,0.9,0.9}{
            \makecell[{{p{\linewidth}}}]{
                \texttt{\tiny{[P1$\langle$GM]}}
                \texttt{guess\_feedback: c<red> r<yellow> a<red> n<yellow> e<yellow>} \\
\\ 
\texttt{Provide your response only in this format.} \\
\texttt{explanation: details} \\
\texttt{guess: word} \\
            }
        }
    }
     \\ \\

    \theutterance \stepcounter{utterance}  
    & \multicolumn{4}{p{0.6\linewidth}}{
        \cellcolor[rgb]{0.561,0.737,0.561}{
            \makecell[{{p{\linewidth}}}]{
                \texttt{\tiny{[P1$\rangle$GM]}}
                \texttt{explanation: The feedback from "crane" (c<red> r<yellow> a<red> n<yellow> e<yellow>) indicates that 'r', 'n', and 'e' are in the word, but not in the 2nd, 4th, and 5th positions respectively. Since 'a' is red, I must avoid it. I need a 5{-}letter word containing r, n, e where 'e' is not at the end, 'r' is not at position 2, and 'n' is not at position 4. "sneer" fits well: it contains s, n, e, e, r. In "sneer", 'n' is at position 2 (which contradicts the yellow 'n' at position 4 in "crane" only if the target has 'n' elsewhere, but wait: in "crane", 'n' was at pos 4 and got yellow, meaning 'n' is in the word but not at pos 4. In "sneer", 'n' is at pos 2. If the target is "sneer", then 'n' is at pos 2. In "crane", 'n' was at pos 4. Since 'n' is yellow in "crane", it means 'n' is in the word but not at pos 4. This is consistent with "sneer" where 'n' is at pos 2. Also 'r' in "crane" was at pos 2 (yellow), so 'r' is not at pos 2. In "sneer", 'r' is at pos 5. This is consistent. 'e' in "crane" was at pos 5 (yellow), so 'e' is not at pos 5. In "sneer", 'e' is at pos 2 and 4. This is consistent. So "sneer" is a strong candidate.} \\
\texttt{guess: sneer} \\
            }
        }
    }
    & & \\ \\

    \theutterance \stepcounter{utterance}  
    & & & \multicolumn{2}{p{0.3\linewidth}}{
        \cellcolor[rgb]{0.9,0.9,0.9}{
            \makecell[{{p{\linewidth}}}]{
                \texttt{\tiny{[GM$|$GM]}}
                \texttt{attempts = 3} \\
\texttt{target\_word = owner} \\
\texttt{guess = sneer} \\
\texttt{guess\_feedback = s<red> n<yellow> e<red> e<green> r<green>} \\
            }
        }
    }
    & & \\ \\

    \theutterance \stepcounter{utterance}  
    & & & \multicolumn{4}{p{0.6\linewidth}}{
        \cellcolor[rgb]{0.9,0.9,0.9}{
            \makecell[{{p{\linewidth}}}]{
                \texttt{\tiny{[P1$\langle$GM]}}
                \texttt{guess\_feedback: s<red> n<yellow> e<red> e<green> r<green>} \\
\\ 
\texttt{Provide your response only in this format.} \\
\texttt{explanation: details} \\
\texttt{guess: word} \\
            }
        }
    }
     \\ \\

    \theutterance \stepcounter{utterance}  
    & \multicolumn{4}{p{0.6\linewidth}}{
        \cellcolor[rgb]{0.561,0.737,0.561}{
            \makecell[{{p{\linewidth}}}]{
                \texttt{\tiny{[P1$\rangle$GM]}}
                \texttt{explanation: The feedback s<red> n<yellow> e<red> e<green> r<green> confirms the word ends in "eer" (positions 4 and 5 are 'e' and 'r'). The 'n' is present but not at position 2, and the 'e' is not at position 3. Since the word ends in "eer", position 4 is 'e' and position 5 is 'r'. The 'n' must be at position 1 or 3. If 'n' is at position 3, the word is "n\_eer" or "nner" (invalid). A common word fitting "n\_eer" is "neener" (too long) or "sneer" (already tried). Wait, if 'n' is at pos 3, the pattern is ? n e e r. But 'e' at pos 3 was red in "sneer", so 'e' cannot be at pos 3. Therefore, 'n' cannot be at pos 3. Thus, 'n' must be at position 1. The pattern is n \_ e e r. The missing letter is at position 2. Common words fitting "n\_eer" include "neer" (not 5 letters), "neener" (6 letters). Let's re{-}evaluate. The word ends in "eer". 'n' is in the word. 'e' is not at pos 3. 'n' is not at pos 2. So 'n' is at pos 1 or 3. If 'n' is at pos 3, the word is ? n e e r. But pos 3 is 'e' in "sneer" which was red. So pos 3 cannot be 'e'. This implies the target word does NOT have 'e' at pos 3. If the target is "n\_eer", then pos 3 is 'e'. Contradiction. Therefore, 'n' cannot be at pos 3. So 'n' must be at pos 1. The word is n \_ e e r. Is there a word like "neer"? No. Maybe "neer" is not the right suffix? Let's re{-}read the feedback for "sneer": s<red> n<yellow> e<red> e<green> r<green>. Pos 4 is 'e', Pos 5 is 'r'. So it ends in "eer". Pos 3 is 'e' in the guess, which is red. So the target does not have 'e' at pos 3. Pos 2 is 'n' in the guess, which is yellow. So 'n' is in the word but not at pos 2. Pos 1 is 's', red. So 's' is out.} \\
            }
        }
    }
    & & \\ \\

    \theutterance \stepcounter{utterance}  
    & & & \multicolumn{2}{p{0.3\linewidth}}{
        \cellcolor[rgb]{0.9,0.9,0.9}{
            \makecell[{{p{\linewidth}}}]{
                \texttt{\tiny{[GM$|$GM]}}
                \texttt{Error: The response should contain the 'guess:' keyword exactly once.} \\
            }
        }
    }
    & & \\ \\

    \theutterance \stepcounter{utterance}  
    & & & \multicolumn{2}{p{0.3\linewidth}}{
        \cellcolor[rgb]{0.9,0.9,0.9}{
            \makecell[{{p{\linewidth}}}]{
                \texttt{\tiny{[GM$|$GM]}}
                \texttt{game\_result = ABORT} \\
            }
        }
    }
    & & \\ \\

\end{supertabular}
}

\end{document}
